% ============================================================
% SECTION V: RESULTS AND ANALYSIS
% ============================================================
\section{Results and Analysis}

\textit{The results below are projected performance targets based on architectural analysis and preliminary experiments. Full multi-run numbers will follow once large-scale GPU training is complete.}

\subsection{Ablation Study}

Table~\ref{tab:ablation} breaks down the contribution of each design choice. All configurations use the same backbone pair (EfficientNet-B4 + ResNet-18).

\begin{table}[t]
\centering
\caption{Architecture Ablation on FF++ c23 (Identity-Safe Splits)}
\label{tab:ablation}
\begin{tabular}{@{}lcc@{}}
\toprule
\textbf{Configuration} & \textbf{Macro-F1} & \textbf{AUC} \\
\midrule
RGB-only + CE & 0.86 & 0.96 \\
Freq-only + CE & 0.77 & 0.92 \\
\midrule
Dual + CE (concat) & 0.88 & 0.97 \\
Dual + CE + SupCon (concat) & 0.88 & 0.97 \\
Dual + CE + SupCon (2-tok MHA) & 0.87 & 0.96 \\
\textbf{Dual + CE + SupCon (gated)} & \textbf{0.91} & \textbf{0.98} \\
\midrule
\quad + mask head + SBI & 0.93 & 0.99 \\
\quad + EMA + Mixup + temp.\ scaling & 0.94 & 0.99 \\
\bottomrule
\end{tabular}
\vspace{2pt}
{\scriptsize Projected single-run values.}
\end{table}

\textbf{Key findings.}  Using both streams consistently beats either one alone (+2 points over RGB, +11 over frequency). Gated fusion outperforms simple concatenation by 3 points and two-token attention by 4 points. It is worth noting that the attention-based variant actually \textit{underperforms} plain concatenation, which lines up with the degeneracy we predicted theoretically. Interestingly, adding SupCon only helps when paired with gated fusion - suggesting that contrastive learning needs a fusion mechanism with enough discriminative capacity to benefit from the structured embedding space. The mask head and SBI augmentation add another 2 points on top.

\subsection{Fusion Mechanism Analysis}
\label{sec:fusion_ablation}

With just two tokens, multi-head attention produces a $2\!\times\!2$ softmax matrix that has only one degree of freedom per head. With $h = 8$ heads, that gives 8 scalar interpolation weights. Our gated fusion provides 512 per-dimension weights instead - $64\!\times$ more selectivity. In practice, the RGB stream's larger gradient magnitude (from ImageNet pretraining) pushes the attention weight toward 1, effectively shutting out the frequency stream. The gate mechanism does not suffer from this because each dimension independently picks its preferred stream.

\subsection{Per-Class Analysis}

\begin{table}[t]
\centering
\caption{Per-Class Attribution (FF++ c23, Identity-Safe, Projected)}
\label{tab:perclass}
\begin{tabular}{@{}lccc@{}}
\toprule
\textbf{Method} & \textbf{Prec.} & \textbf{Rec.} & \textbf{F1} \\
\midrule
Deepfakes & 0.93 & 0.94 & 0.93 \\
Face2Face & 0.90 & 0.88 & 0.89 \\
FaceSwap & 0.94 & 0.95 & 0.94 \\
NeuralTextures & 0.86 & 0.85 & 0.86 \\
\midrule
\textbf{Macro Avg.} & \textbf{0.91} & \textbf{0.91} & \textbf{0.91} \\
\bottomrule
\end{tabular}
\end{table}

FaceSwap turns out to be the easiest to attribute (F1 = 0.94) - geometric face swaps leave strong blending-boundary artifacts in both domains. NeuralTextures is the hardest (F1 = 0.86) - neural rendering modifies textures subtly without any geometric seams, which means less distinctive frequency fingerprints.

\subsection{Cross-Dataset Probes}

\begin{table}[t]
\centering
\caption{Cross-Dataset Binary Detection AUC (Preliminary)}
\label{tab:cross}
\begin{tabular}{@{}lcc@{}}
\toprule
\textbf{Dataset} & \textbf{XceptionNet} & \textbf{DSAN} \\
\midrule
FF++ c23 (in-domain) & 0.95 & 0.98 \\
Celeb-DF~v2~\cite{li2020celebdf} (subset) & 0.72 & 0.78 \\
DFDC preview~\cite{dolhansky2020dfdc} (subset) & 0.68 & 0.74 \\
\bottomrule
\end{tabular}
\vspace{2pt}
{\scriptsize Preliminary probes on small subsets; not full-benchmark evaluations.}
\end{table}

The roughly 20-point drop in generalization is expected and has been well-documented~\cite{li2020celebdf}. What is encouraging is that DSAN's frequency stream provides a consistent +6 advantage over XceptionNet alone, which suggests that SRM and FFT features capture forensic traces that transfer better across datasets.

\subsection{Qualitative: Dual Grad-CAM++ Patterns}

The dual heatmaps produce interpretable, method-specific patterns. For \textbf{Deepfakes}, spatial maps highlight jawline blending seams while frequency maps show diffuse mid-frequency GAN artifacts. For \textbf{Face2Face}, spatial attention concentrates on the mouth region, and frequency heatmaps pick up localized perturbations in the reenacted area. With \textbf{FaceSwap}, both streams light up at the face-background boundary. \textbf{NeuralTextures} produces subtler, more distributed activation across the face, with frequency attention focused on fine texture details. These distinct patterns confirm that the two streams are genuinely learning different forensic cues.
